% 連続体力学 連続体力学-4 問題1（027-volume） 解いた日 2026-11-25
%% kotae
\toi{1}{△}{記述}
\begin{align*}
\det(\tens{I}+\tens{D})&=\Bigl(1+\pdv{r_{1}}{x_{1}}\Bigr)\Bigl(1+\pdv{r_{2}}{x_{2}}\Bigr)\\
&\qquad\times\Bigl(1+\pdv{r_{3}}{x_{3}}\Bigr)\nuke{+(\text{2次以上})}\\
&\ayamari{=}{\simeq}1+\underbrace{\pdv{r_{1}}{x_{1}}+\pdv{r_{2}}{x_{2}}+\pdv{r_{3}}{x_{3}}}_{\Tr(\tens{E})}
\end{align*}
\[ \therefore\ \frac{\delta V'}{\delta V}=1+\Tr(\tens{E}) \]
\[ \therefore\ \frac{\delta V'-\delta V}{\delta V}=\Tr(\tens{E}) \]
\naze{行列式は置換についての和で，対角成分の積は恒等置換の項にすぎない。ほかの置換の項は非対角成分 $D_{i\sigma(i)}$ を2つ以上含むので $\tens{D}$ の2次以上として捨てられる，という一言が要の論点である（答案に必要な論点）。これが無いと1行目は一般には誤った等式になる。恒等置換の項も展開して2次以上を捨てるので2行目は $\simeq$。また $\partial r_{i}/\partial x_{i}=D_{ii}=E_{ii}$ となるのは $\tens{F}$ の対角成分が0だからである。}
\toi{2}{△}{計算}
\[ \tens{D}=\begin{pmatrix}a&\gamma&0\\0&b&0\\0&0&c\end{pmatrix},\quad \tens{E}=\begin{pmatrix}a&\gamma/2&0\\\gamma/2&b&0\\0&0&c\end{pmatrix} \]
\[ \tens{F}=\begin{pmatrix}\ayamari{a}{0}&\gamma/2&0\\-\gamma/2&\ayamari{b}{0}&0\\0&0&\ayamari{c}{0}\end{pmatrix} \]
\naze{$\tens{F}=(\tens{D}-\tens{D}^{\mathrm{T}})/2$ は反対称で，対角成分は必ず0である。対角の $a,b,c$ は $\tens{E}$ だけに入る（書いたままだと $\tens{E}+\tens{F}$ の対角が $2a,2b,2c$ になり $\tens{D}$ に戻らない）。}
\[ \Tr(\tens{E})=a+b+c \]
\toi{3}{◎}{}
\[ \det\begin{pmatrix}1+a&\gamma&0\\0&1+b&0\\0&0&1+c\end{pmatrix}=(1+a)(1+b)(1+c) \]
体積変化率
\[ (1+a)(1+b)(1+c)-1\simeq\num{1.498499e-3} \]
\[ a+b+c=\num{1.50e-3} \]
\[ \text{差}=\num{-1.5e-6} \]
（歪みの2次以上の量）
%% kizuki
\item 行列式の1次の展開は「置換についての和」から考える。恒等置換以外の項は非対角成分を2つ以上含むので2次以上になり，残るのは対角成分の和だけである。「$\det$＝対角成分の積」が厳密に成り立つのは (3) のような三角行列のときだけなので，(1) では捨てた項が2次以上である理由を一言書く。
\item 反対称テンソル $\tens{F}$ の対角成分は常に0である。$\tens{E}$・$\tens{F}$ を書いたら，$\tens{E}+\tens{F}=\tens{D}$ と $\tens{F}^{\mathrm{T}}=-\tens{F}$ の2つで検算する。
\item \mae{第3回 問題5 (2)・10-28} 「示せ」の問いで，主張の要の段（前回は分解の一意性，今回は1次の項だけが残る理由）を飛ばした。示す式の各等号が一般に成り立つかを見直し，近似で捨てた項には「2次以上」と書き添える。
